%Das ist die Hauptdatei. 
%Diese Datei ist die, welche von pdflatex übersetzt wird (erzeugt dann pdf). 
%Dieses Liederbuch nutzt das Latex Songs Package, welches unter der GPLv2 steht. 
%Eine Kopie der GPLv2 kann hier eingesehen werden: https://opensource.org/licenses/gpl-2.0.php 
%Kurzum sagt die Lizenz, dass das Paket genutzt werden darf, und dass der Entwickler des Paketes 
%für nichts haftbar gemacht werden kann. 
%Durch die Nutzung des Paketes steht dieses Liederbuch automatisch auch unter der GPLv2 Lizenz. 
%Für uns heißt das, wenn wir nach Quellcode gefragt werden, diesen auch ausgeben müssen. 
%Wer unser Liederbuch modifiziert, muss bei Nachfrage ebenfalls sein Buch offenlegen. 
%Es muss kenntlich sein, dass Nutzer des Liederbuches diese Rechte haben, 
%somit muss (am besten in diesem Buch selbst) auf die Lizenz hingewiesen werden. 
%Ein entsprechender Verweis wurde auf der Titelseite angebracht. 
%Ein weiteres Problem stellt das Urheberrecht dar, wenn man ein solches Liederbuch bastelt. 
%Tatsächlich ist es so, dass wir den gesamten Quellcode rausrücken müssen, wenn wir gefragt werden. 
%Wir dürfen aber die vom Urheberrecht geschützen Liedtexte nicht einfach öffentlich zum Download anbieten. 
%Mein Ansatz hierfür wäre, die Lieder als separates Projekt anzusehen und auch so zu pflegen,
%denn dann trifft die GPLv2 nicht auf die Lieder sondern nur auf das Master-Layout zu. 
%Die offizielle Version des letzten Satzes ist dann: "Wir nutzen nur Urheberrechtlich ungeschütztes Material 
%und solches, von welchem wir eine Genehmigung haben"

\documentclass{book}
\usepackage[lmargin=10mm, rmargin=5mm, tmargin=2mm, bmargin=2mm, text={15cm,10cm},centering]{geometry}
\geometry{papersize={148mm,105mm}}
\usepackage[bookmarks]{hyperref}
\usepackage[chorded]{songs}
\usepackage{graphicx}
\usepackage[ngerman]{babel}

\pagestyle{empty}

\newindex{titleidx}{cbtitle}
\newauthorindex{authidx}{cbauth}
\newscripindex{scripidx}{cbscrip}

\songcolumns{1}
\usepackage[parfill]{parskip}
\renewcommand{\lyricfont}{\sffamily\small}
\renewcommand\printchord[1]{\sffamily\textmd\large#1}
\renewcommand{\stitlefont}{\sffamily\textmd}
\versesep=8pt

\begin{document}
\begin{titlepage}
\vspace*{\stretch{1.0}}
\begin{center}
{\sffamily\huge\bfseries Klampfenkrampf\par}
{\sffamily\small\bfseries Liederbuch\par}
\vspace{0.5cm}
{\sffamily\Large DPSG Neckarelz-Diedesheim\par}
\vspace{0.5cm}
{\sffamily\Large Stamm Don Bosco\par}
\vspace{0.5cm}
\begin{figure}[h]
\centering
\includegraphics[width=0.4\textwidth]{img/logo.jpg}\par\
\end{figure}
{\sffamily\small Zuletzt am \today{} aktualisiert\par}
\vspace{0.5cm}
{\sffamily\tiny Das Layout dieses Buches steht unter der Gnu General Public License Version 2.0.\par}
\end{center}
\vspace*{\stretch{2.0}}
\end{titlepage}


\begin{songs}{titleidx,authidx,scripidx} 
	\beginsong{Lagerboogie}[by={}]
\beginverse
Wir kom\[G]men aus dem Häusermeer und fahren in die \[D]Welt, Wir ziehen immer kreuz und quer, so lang es uns gef\[G]ällt.
\endverse

\beginchorus
Ja ja ja, ts\[C]chu tschu der Lagerboogie \[G]ist unser Boogie-Woogie, ts\[D]chu tschu tschu, die Zeit vergeht im N\[G]u!    (2x)
\endchorus

\beginverse
Kolum\[G]bus hat die Welt entdeckt und die ist kugelr\[D]und, es sah ein Knab ein Rösslein stehn mit 180 Pf\[G]und.
\endverse

\beginverse
Die K\[G]uh gibt Süß- und Sauermilch den lieben langen T\[D]ag, der Ochse, dieses dumme Vieh, der gibt uns nur Spi\[G]nat.
\endverse

\beginverse
Und we\[G]nn du denkst, du hast ihn schon, den goldnen Abends\[D]tern, dann kriegst de eins mim Pann vorn Kopf, das ist der Tag des He\[G]rrn.
\endverse

\beginverse
Die L\[G]agerruhe ist um zehn, wir sagen gute N\[D]acht, wir müssen in die Falle gehn, die Lagerleitung w\[G]acht.
\endverse

\vfill

\hfill \gtab{C}{X32010} \hfill \gtab{D}{XX0232} \hfill \gtab{G}{320033} \hspace*{\fill}

\endsong



\end{songs}

\showindex{Inhaltsverzeichnis}{titleidx}
%nicht genutzte indizes für andere inhaltsverzeichnisse
%sollen diese genutzt werden, muss das makefile angepasst werden!
%\showindex{Index of Authors and Composers}{authidx}
%\showindex{Index of Scripture}{scripidx}

\end{document}

